\begin{center}\textit{Ein neuer Algorithmus zur Entscheidung, ob eine gegebene positive ganze Zahl quadratischer Rest oder Nichtrest einer gegebenen positiven Primzahl ist.}\end{center}

\Article{1}

Bevor wir die neue Lösung dieser Aufgabe darstellen, wollen wir die in den \textit{Disquisitiones Arithmeticae} angegebene Lösung kurz wiederholen. Sie lässt sich mit Hilfe des Fundamentaltheorems und der folgenden bekannten Sätze recht rasch durchführen:

I. Die Beziehung der Zahl $a$ zur Zahl $b$ (ob jene ein quadratischer Rest oder Nichtrest dieser ist) ist dieselbe wie die der Zahl $c$ zu $b$, wenn $a\equiv c\ (\mathrm{mod.}\ b)$ gilt.

II. Ist $a$ das Produkt der Faktoren $\alpha,\beta,\gamma,\delta$ usw. und $b$ eine Primzahl, so hängt die Beziehung von $a$ zu $b$ folgendermaßen von der Beziehung dieser Faktoren zu $b$ ab: $a$ ist ein quadratischer Rest von $b$ oder ein Nichtrest, je nachdem, ob eine gerade oder ungerade Anzahl dieser Faktoren Nichtreste von $b$ sind. Ein Faktor, der ein Quadrat ist, braucht daher bei dieser Untersuchung überhaupt nicht berücksichtigt zu werden; ist ein Faktor dagegen eine Potenz einer ganzen Zahl mit ungeradem Exponenten, so kann diese ganze Zahl selbst an seine Stelle treten.

III. Die Zahl 2 ist ein quadratischer Rest jeder Primzahl der Form $8m+1$ oder $8m+7$, dagegen ein Nichtrest jeder Primzahl der Form $8m+3$ oder $8m+5$.

\SourceLeaf{70}%
Es sei also eine Zahl $a$ gegeben, deren Beziehung zu einer Primzahl $b$ gesucht wird. Für $a$ ist, falls sie größer als $b$ ist, zunächst ihr kleinster positiver Rest modulo $b$ einzusetzen. Zerlegt man diesen Rest in seine Primfaktoren, so führt Satz II die Frage darauf zurück, die Beziehung jedes dieser Faktoren zu $b$ zu bestimmen. Die Beziehung des Faktors $2$ (wenn er einmal, dreimal, fünfmal usw. vorkommt) ist nach Satz III bekannt; die Beziehung jedes übrigen Faktors hängt nach dem Fundamentaltheorem von der Beziehung von $b$ selbst zu diesem Faktor ab. Anstelle der einen Beziehung der gegebenen Zahl zur Primzahl $b$ sind also nun mehrere Beziehungen der Zahl $b$ zu anderen ungeraden Primzahlen zu untersuchen, die kleiner als $b$ sind. Diese Aufgaben werden auf dieselbe Weise auf kleinere Moduln zurückgeführt, und diese aufeinanderfolgenden Reduktionen müssen offenbar schließlich enden.

\Article{2}

Zur Erläuterung dieser Lösung sei die Beziehung der Zahl $103$ zu $379$ gesucht. Da $103$ bereits kleiner als $379$ und selbst eine Primzahl ist, ist unmittelbar das Fundamentaltheorem anzuwenden. Es besagt, dass die gesuchte Beziehung derjenigen der Zahl $379$ zu $103$ entgegengesetzt ist. Diese ist wiederum gleich der Beziehung der Zahl $70$ zu $103$, die ihrerseits von den Beziehungen der Zahlen $2,\ 5,\ 7$ zu $103$ abhängt. Die erste dieser Beziehungen ist nach Satz III bekannt. Die zweite hängt nach dem Fundamentaltheorem von der Beziehung der Zahl $103$ zu $5$ ab, die nach Satz I gleich der Beziehung von $3$ zu $5$ ist; diese hängt wiederum nach dem Fundamentaltheorem von der Beziehung von $5$ zu $3$ ab, die nach Satz I gleich der nach Satz III bekannten Beziehung von $2$ zu $3$ ist. Ebenso hängt die Beziehung von $7$ zu $103$ nach dem Fundamentaltheorem von der Beziehung von $103$ zu $7$ ab, die nach Satz I gleich der Beziehung von $5$ zu $7$ ist; diese hängt wiederum nach dem Fundamentaltheorem von der Beziehung von $7$ zu $5$ ab, die nach Satz I gleich der nach Satz III bekannten Beziehung von $2$ zu $5$ ist. Wollen wir nun diese Analyse in eine Synthese umsetzen, so lässt sich die Entscheidung der Frage auf vierzehn Schritte zurückführen. Wir geben sie hier vollständig an, damit die größere Eleganz der neuen Lösung umso deutlicher hervortritt.

{\raggedright\hangindent=2.2em\hangafter=1\noindent 1.\quad Die Zahl $2$ ist ein quadratischer Rest der Zahl $103$ \mbox{(Satz III)}.\par}

{\raggedright\hangindent=2.2em\hangafter=1\noindent 2.\quad Die Zahl $2$ ist ein quadratischer Nichtrest der Zahl $3$ \mbox{(Satz III)}.\par}

{\raggedright\hangindent=2.2em\hangafter=1\noindent 3.\quad Die Zahl $5$ ist ein quadratischer Nichtrest der Zahl $3$ \mbox{(aus I und 2)}.\par}

{\raggedright\hangindent=2.2em\hangafter=1\noindent 4.\quad Die Zahl $3$ ist ein quadratischer Nichtrest der Zahl $5$ \mbox{(Fundamentaltheorem und 3)}.\par}

{\raggedright\hangindent=2.2em\hangafter=1\noindent 5.\quad Die Zahl $103$ ist ein quadratischer Nichtrest der Zahl $5$ \mbox{(I und 4)}.\par}

\SourceLeaf{71}%
{\raggedright\hangindent=2.2em\hangafter=1\noindent 6.\quad Die Zahl $5$ ist ein quadratischer Nichtrest der Zahl $103$ \mbox{(Fundamentaltheorem und 5)}.\par}

{\raggedright\hangindent=2.2em\hangafter=1\noindent 7.\quad Die Zahl $2$ ist ein quadratischer Nichtrest der Zahl $5$ \mbox{(Satz III)}.\par}

{\raggedright\hangindent=2.2em\hangafter=1\noindent 8.\quad Die Zahl $7$ ist ein quadratischer Nichtrest der Zahl $5$ \mbox{(I und 7)}.\par}

{\raggedright\hangindent=2.2em\hangafter=1\noindent 9.\quad Die Zahl $5$ ist ein quadratischer Nichtrest der Zahl $7$ \mbox{(Fundamentaltheorem und 8)}.\par}

{\raggedright\hangindent=2.2em\hangafter=1\noindent 10.\quad Die Zahl $103$ ist ein quadratischer Nichtrest der Zahl $7$ \mbox{(I und 9)}.\par}

{\raggedright\hangindent=2.2em\hangafter=1\noindent 11.\quad Die Zahl $7$ ist ein quadratischer Rest der Zahl $103$ \mbox{(Fundamentaltheorem und 10)}.\par}

{\raggedright\hangindent=2.2em\hangafter=1\noindent 12.\quad Die Zahl $70$ ist ein quadratischer Nichtrest der Zahl $103$ \mbox{(II, 1, 6, 11)}.\par}

{\raggedright\hangindent=2.2em\hangafter=1\noindent 13.\quad Die Zahl $379$ ist ein quadratischer Nichtrest der Zahl $103$ \mbox{(I und 12)}.\par}

{\raggedright\hangindent=2.2em\hangafter=1\noindent 14.\quad Die Zahl $103$ ist ein quadratischer Rest der Zahl $379$ \mbox{(Fundamentaltheorem und 13)}.\par}

Im Folgenden verwenden wir der Kürze wegen das in \textit{Comment. Gotting. Vol. XVI} eingeführte Zeichen. Mit $[x]$ bezeichnen wir die Größe $x$ selbst, wenn $x$ eine ganze Zahl ist, und die größte ganze Zahl unterhalb von $x$, wenn $x$ nicht ganzzahlig ist. Damit ist $x-[x]$ stets nichtnegativ und kleiner als eins.

\Article{3}

\textsc{Aufgabe.}\quad \textit{Es seien $a,\ b$ positive, zueinander teilerfremde ganze Zahlen, und es sei $[\frac12a]=a'$ gesetzt. Gesucht ist die Summe}

\[
\gbr{\frac ba}+\gbr{\frac{2b}a}+\gbr{\frac{3b}a}+\gbr{\frac{4b}a}+\text{usw.}+\gbr{\frac{a'b}a}
\]

\textsc{Lösung.}\quad Bezeichnen wir eine solche Summe der Kürze wegen mit $\varphi(a,b)$, so gilt auch

\[
\varphi(b,a)=\gbr{\frac ab}+\gbr{\frac{2a}b}+\gbr{\frac{3a}b}+\text{usw.}+\gbr{\frac{b'a}b}
\]

wenn wir $[\frac12b]=b'$ setzen. Im dritten Beweis des Fundamentaltheorems wurde gezeigt, dass für ungerade $a$ und $b$ gilt:

\[
\varphi(a,b)+\varphi(b,a)=a'b'
\]

Mit derselben Methode lässt sich die Gültigkeit dieses Satzes leicht auch auf den Fall ausdehnen, dass eine der Zahlen $a,\ b$ ungerade ist, wie wir dort bereits angedeutet haben. Man teile $a$ durch $b$, wie bei dem Verfahren zur Bestimmung des größten gemeinsamen Teilers zweier ganzer Zahlen; dabei sei $\beta$ der Quotient und $c$ der Rest. Dann teile man $b$ durch $c$ usw., so dass die folgenden Gleichungen entstehen:

\SourceLeaf{72}%
\[
\begin{aligned}
a&=\beta b+c\\
b&=\gamma c+d\\
c&=\delta d+e\\
d&=\varepsilon e+f\ \text{usw.}
\end{aligned}
\]

Auf diese Weise erreichen wir in der ständig abnehmenden Folge der Zahlen $b,c,d,e,f$ usw. schließlich die Eins, da $a$ und $b$ nach Voraussetzung teilerfremd sind. Die letzte Gleichung lautet daher

\[
k=\lambda l+1
\]

Da offenbar

\[
\begin{aligned}
\gbr{\frac ab}&=\gbr{\beta+\frac cb}=\beta+\gbr{\frac cb}\\
\gbr{\frac{2a}b}&=\gbr{2\beta+\frac{2c}b}=2\beta+\gbr{\frac{2c}b}\\
\gbr{\frac{3a}b}&=\gbr{3\beta+\frac{3c}b}=3\beta+\gbr{\frac{3c}b}
\end{aligned}
\]

usw.. gilt, erhalten wir

\[
\varphi(b,a)=\varphi(b,c)+\tfrac12\beta(b'b'+b')
\]

und damit

\[
\varphi(a,b)=a'b'-\tfrac12\beta(b'b'+b')-\varphi(b,c)
\]

\Needspace{10\baselineskip}
Setzen wir $[\frac12c]=c',\ [\frac12d]=d'\ [\frac12e]=e'$ usw., so ergibt sich durch ähnliche Überlegungen

\[
\begin{aligned}
\varphi(b,c)&=b'c'-\tfrac12\gamma(c'c'+c')-\varphi(c,d)\\
\varphi(c,d)&=c'd'-\tfrac12\delta(d'd'+d')-\varphi(d,e)\mathord{.}\\
\varphi(d,e)&=d'e'-\tfrac12\varepsilon(e'e'+e')-\varphi(e,f)
\end{aligned}
\]

usw., bis hin zu

\[
\varphi(k,l)=k'l'-\tfrac12\lambda(l'l'+l')-\varphi(l,1)
\]

Da offenbar $\varphi(l,1)=0$ ist, erhalten wir hieraus die Formel

\[
\begin{aligned}
\varphi(a,b)&=a'b'-b'c'+c'd'-d'e'+\text{usw.}\pm k'l'\\
&\quad-\tfrac12\beta(b'b'+b')+\tfrac12\gamma(c'c'+c')-\tfrac12\delta(d'd'+d')\\
&\quad+\tfrac12\varepsilon(e'e'+e')-\text{usw.}\mp\tfrac12\lambda(l'l'+l')
\end{aligned}
\]

\Article{4}

Aus dem dritten Beweis folgt nun leicht, dass die Beziehung der Zahl $b$ zu $a$, sofern $a$ eine Primzahl ist, unmittelbar am Wert der Summe\SourceLeaf{73} $\varphi(a,2b)$ erkannt wird. Je nachdem, ob diese Summe gerade oder ungerade ist, ist $b$ ein quadratischer Rest von $a$ oder ein Nichtrest. Zu demselben Zweck kann auch die Summe $\varphi(a,b)$ selbst verwendet werden, allerdings mit der Einschränkung, dass der Fall eines ungeraden $b$ von dem eines geraden zu unterscheiden ist. Es gilt nämlich:

I.\quad Ist $b$ ungerade, so ist $b$ ein quadratischer Rest oder Nichtrest von $a$, je nachdem, ob $\varphi(a,b)$ gerade oder ungerade ist.

II.\quad Ist $b$ gerade, so gilt dieselbe Regel, wenn außerdem $a$ die Form $8n+1$ oder $8n+7$ hat; hat dagegen bei einem geraden Wert von $b$ der Modul $a$ die Form $8n+3$ oder $8n+5$, so ist die entgegengesetzte Regel anzuwenden: $b$ ist ein quadratischer Rest von $a$, wenn $\varphi(a,b)$ ungerade ist, dagegen ein Nichtrest, wenn $\varphi(a,b)$ gerade ist.

All dies folgt sehr leicht aus Artikel 4 des dritten Beweises.

\newpage
\Article{5}

\textit{Beispiel.} Ist die Beziehung der Zahl $103$ zur Primzahl $379$ gesucht, so haben wir zur Berechnung der Summe $\varphi(379,103)$:

\[
\begin{array}{r|r|l}
a=379&a'=189&\\
b=103&b'=51&\beta=3\\
c=70&c'=35&\gamma=1\\
d=33&d'=16&\delta=2\\
e=4&e'=2&\varepsilon=8
\end{array}
\]

und damit

\[
\begin{aligned}
\varphi(379,103)&=9639-1785+560-32-3978+630-272+24\\
&=4786
\end{aligned}
\]

Also ist $103$ ein quadratischer Rest der Zahl $379$. Wollen wir stattdessen die Summe $(379,206)$ für denselben Zweck verwenden, so erhalten wir das folgende Schema:

\[
\begin{array}{r|r|l}
379&189&\\
206&103&1\\
173&86&1\\
33&16&5\\
8&4&4
\end{array}
\]

Hieraus ergibt sich

\SourceLeaf{74}%
\[
\begin{aligned}
\varphi(379,206)&=19467-8858+1376-64-5356+3741-680+40\\
&=9666
\end{aligned}
\]

Somit ist $103$ ein quadratischer Rest der Zahl $379$.

\Needspace{28\baselineskip}
\Article{6}

Um die Beziehung der Zahl $b$ zu $a$ zu entscheiden, brauchen die einzelnen Glieder der Summe $\varphi(a,b)$ nicht berechnet zu werden; es genügt zu wissen, wie viele von ihnen ungerade sind. Unsere Regel lässt sich daher auch so angeben:

Es sei wie oben $a=\beta b+c,\ b=\gamma c+d,\ c=\delta d+e$ usw., bis in der Folge der Zahlen $a,\ b,\ c,\ d,\ e$ usw. die Eins erreicht ist. Man setze $[\frac12a]=a',\ [\frac12b]=b',\ [\frac12c]=c'$ usw. Ferner sei $\mu$ die Anzahl der ungeraden Zahlen in der Folge $a',b',c'$ usw., auf die unmittelbar eine ungerade Zahl folgt, und $\nu$ die Anzahl der ungeraden Zahlen in der Folge $\beta,\gamma,\delta$ usw., denen in der Folge $b',c',d'$ usw. jeweils eine Zahl der Form $4n+1$ oder $4n+2$ entspricht. Dann ist $b$ ein quadratischer Rest oder Nichtrest von $a$, je nachdem, ob $\mu+\nu$ gerade oder ungerade ist. Ausgenommen ist allein der Fall, dass zugleich $b$ gerade ist und $a$ die Form $8n+3$ oder $8n+5$ hat; hier gilt die entgegengesetzte Regel.

In unserem Beispiel folgt in der Folge $a',b',c',d',e'$ zweimal unmittelbar auf eine ungerade Zahl wieder eine ungerade; also ist $\mu=2$; in der Folge $\beta',\gamma',\delta',\varepsilon'$ kommen zwar zwei ungerade Zahlen vor, ihnen entsprechen jedoch in der Folge $b',c',d',e'$ Zahlen der Form $4n+3$, so dass $\nu=0$ ist. Damit ist $\mu+\nu$ gerade, und folglich ist $103$ ein quadratischer Rest der Zahl $379$.

\BookRule{42mm}
